% !PRINTED 267
\weberpage{267}
%<<CONT
reell wird. Es ergiebt sich aber nach (1) und (2)
\[
R\,S\,R^{-1} = \begin{pmatrix} \lambda\,\varrho\,\sigma - \mu\,\nu\,\sigma^{-1}, & -\lambda\,\mu\,(\sigma-\sigma^{-1}) \\ \varrho\,\nu\,(\sigma-\sigma^{-1}), & -\mu\,\nu\,\sigma+\varrho\,\lambda\,\sigma^{-1} \end{pmatrix},
\]
und dies wird reell, wenn $\lambda\mu$ und $\varrho\nu$ rein imaginär, $\lambda\varrho$ und
$-\mu\nu$ conjugirt imaginär sind. Diesen Bedingungen kann man
auf mehrfache Art genügen: am einfachsten wohl, und zwar
zugleich so, dass die Determinante von $R = 1$ wird, wenn man
\weq{4}{R = \begin{pmatrix} \frac{1}{2}, & -\frac{1}{2}\,\varepsilon^{-1} \\ \varepsilon, & 1 \end{pmatrix}}
setzt. Ist andererseits
\[
A = \begin{pmatrix} a, & b \\ c, & d \end{pmatrix}
\]
eine beliebige reelle Substitution aus $L_p$, so ergiebt sich nach (4)
wegen $\varepsilon^2 = N$:
\[
R^{-1}\,A\,R =
\]
\weq{5}{\begin{pmatrix} \dfrac{a+d}{2}+b\,\varepsilon+\dfrac{c}{4}\,\varepsilon^{-1}, & -\dfrac{a-d}{2}\,\varepsilon^{-1}+b-\dfrac{c}{4}\,N^{-1} \\[1.2ex] -N\left(\dfrac{a-d}{2}\,\varepsilon^{-1}+b-\dfrac{c}{4}\,N^{-1}\right), & \dfrac{a+d}{2}-b\,\varepsilon-\dfrac{c}{4}\,\varepsilon^{-1} \end{pmatrix},}
wofür wir auch
\weq{6}{A = \begin{pmatrix} \alpha, & \beta \\ -N\beta', & \alpha' \end{pmatrix} = \begin{pmatrix} \alpha, & \beta \\ -N\beta^p, & \alpha^p \end{pmatrix}, \quad \alpha\alpha'+N\beta\beta' = 1}
setzen, worin $\alpha,\alpha'$ und $\beta,\beta'$ zwei conjugirt imaginäre Paare
sind, so dass nach §.\,65, (4) $\alpha' = \alpha^p$, $\beta' = \beta^p$ gesetzt werden
kann.

Wenn wir umgekehrt irgend eine Substitution $A$ von der
Form (6) nehmen, so ergiebt sich
\[
A = R\,A\,R^{-1} =
\]
\weq{7}{\begin{pmatrix} \dfrac{\alpha+\alpha'}{2}+\varepsilon\,\dfrac{\beta'-\beta}{2}, & \dfrac{\beta+\beta'}{4}+\varepsilon^{-1}\,\dfrac{\alpha-\alpha'}{4} \\[1.2ex] -N(\beta+\beta')+\varepsilon(\alpha-\alpha'), & \dfrac{\alpha+\alpha'}{2}-\varepsilon\,\dfrac{\beta'-\beta}{2} \end{pmatrix}}
als eine \gesperrt{reelle} Substitution aus $L_p$.

\gesperrt{Daraus geht hervor, dass, wenn $A$ die ganze Gruppe $L_p$
durchläuft, die durch (5) und (6) bestimmte Substitution $A$ eine
mit $L_p$ isomorphe Gruppe durchläuft, die wir mit $\Lambda_p$ bezeichnen}.
